% !TEX TS-program = pdflatexmk
\documentclass[12pt]{amsart}
\usepackage{multicol}
\usepackage{float}

%\usepackage[parfill]{parskip}    % Activate to begin paragraphs with an empty line rather than an indent

\usepackage[margin=1in]{geometry}


\usepackage{amsmath,amssymb,amsthm,latexsym,graphicx}
\usepackage[normalem]{ulem}
\usepackage{setspace} %used for doublespacing, etc.
\usepackage{hyperref}
\usepackage{cancel}
\usepackage[dvipsnames,usenames]{color}
\usepackage[all]{xy}
\usepackage{fancyhdr}
\pagestyle{fancy}
	\renewcommand{\headrulewidth}{0.5pt} % and the line
	\headsep=1cm
	
\DeclareGraphicsRule{.tif}{png}{.png}{`convert #1 `dirname #1`/`basename #1 .tif`.png}

%Some useful environments.
\newtheorem{theorem}{Theorem}
\newtheorem{corollary}[theorem]{Corollary}
\newtheorem{conjecture}[theorem]{Conjecture}
\newtheorem{lemma}[theorem]{Lemma}
\newtheorem{proposition}[theorem]{Proposition}
\newtheorem{definition}[theorem]{Definition}
\newtheorem{example}[theorem]{Example}
\newtheorem{axiom}{Axiom}
\theoremstyle{remark}
\newtheorem{remark}{Remark}
\newtheorem*{exercise}{Exercise}%[section]

%Some useful shortcuts for our favorite fields
\def\RR{\ensuremath{\mathbb R}} %Note, you can use these WITHOUT entering math  \mod e
\def\NN{\ensuremath{\mathbb N}}
\def\ZZ{\ensuremath{\mathbb Z}}
\def\QQ{{\ensuremath\mathbb Q}}
\def\CC{\ensuremath{\mathbb C}}
\def\EE{{\ensuremath\mathbb E}}

%Some useful shortcuts for formatting lists
\newcommand{\bc}{\begin{center}}
\newcommand{\ec}{\end{center}}
\newcommand{\be}{\begin{enumerate}}
\newcommand{\ee}{\end{enumerate}}
\newcommand{\bi}{\begin{itemize}}
\newcommand{\ei}{\end{itemize}}

%Some useful shortcuts for formatting mathematical symbols
\newcommand{\ol}[1]{\overline{#1}}
\newcommand{\oimp}[1]{\overset{#1}{\iff}} %labeled iff symbol
\newcommand{\bv}[1]{\ensuremath{ \vec{\mathbf{#1}}} } %makes a vector.
\newcommand{\mc}[1]{\ensuremath{\mathcal{#1}}} %put something in caligraphic font
\newcommand{\mpg}[1]{\marginpar{ #1}} %to put comments in margins
\newcommand{\bsl}[1]{\texttt{\symbol{92}{\em #1}}} %for backslashes.
\newcommand{\normale}{\trianglelefteq}
\newcommand{\normal}{\triangleleft}

%Code for formatting the proofs a little nicer for submitted homework
\makeatletter
\renewenvironment{proof}[1][\proofname]{\par\doublespacing
  \pushQED{\qed}%
  \normalfont \topsep6\p@\@plus6\p@\relax
  \list{}{%
    \settowidth{\leftmargin}{\itshape\proofname:\hskip\labelsep}%
    \setlength{\labelwidth}{0pt}%
    \setlength{\itemindent}{-\leftmargin}%
  }%
  \item[\hskip\labelsep\itshape#1\@addpunct{:}]\ignorespaces
}{%
  \popQED\endlist\@endpefalse
  \singlespacing
}
\makeatother


%Commenting tools. You can ignore this stuff unless we're exchanging materials where I'm commenting in detail on how you are writing/using LaTeX.
\usepackage{soul}
\definecolor{highlight}{rgb}{1,0.6,0.6}
\sethlcolor{highlight}
\newcommand{\hlm}[1]{\colorbox{highlight}{$\displaystyle #1$}}
\newtheoremstyle{mycomment}{\topsep}{-0in}{\small \itshape \sffamily}{}{\small \itshape\sffamily}{:}{.5em}{}
\theoremstyle{mycomment}
\newtheorem*{acomment}{\color{BrickRed}{Comment}}
\newcommand{\com}[1]{{\color{OliveGreen}\begin{acomment}{#1} %#2 \color{black} 
\end{acomment}\noindent}}
%\newcommand{\com}[1]{{\color{BrickRed}{\\ Comment:}\color{OliveGreen}{#1} \\}}
\newcommand{\red}[1]{{\color{BrickRed} #1}}
\newcommand{\blue}[1]{{\color{MidnightBlue}#1}}
\newcommand{\green}[1]{{\color{OliveGreen}#1}}
\newcommand{\mwrong}[2]{\red{\cancel{#1}}\green{#2}}
\newcommand{\wrong}[2]{\red{\sout{#1}}\green{#2}}
\definecolor{OliveGreen}{rgb}{.3,.5,.2}
\definecolor{MidnightBlue}{rgb}{.3,.4,.6}
\newcommand{\pts}[1]{\hfill\blue{{#1}/5}}

\chead{MATH 631F}
\pagestyle{fancy}
% \mod ify these items:
\rhead{\emph{Stefano Fochesatto}}
\lhead{\emph{HW \#11 --- \today}}

\DeclareMathOperator{\lcm}{lcm}
\begin{document}

\thispagestyle{fancy}
\author{Coordinator: Coordinator's Name}






%7.4 7, 30, 41
%9.1 1, 4, 9, 11, 13.

\section*{\textbf{Section 7.4}}

\begin{exercise}[\textbf{7.4.7}] Let $R$ be a commutative ring with 1. Prove that the principal ideal generated by $x$
  in the polynomial ring $R[x]$ is a prime ideal if and only if $R$ is an integral domain. Prove that $(x)$ is a 
  maximal ideal if and only if $R$ is a field. 
  \begin{proof} Let $R$ be a commutative ring with 1. Consider the homomorphism $\varphi: R[x] \to R$ defined by 
    $\varphi(a_0 + a_1x + a_2x^2 + \dots) = a_0$. Let $p(x), q(x) \in R[x]$ such that $\varphi(p(x)) = a_0$ and $\varphi(q(x)) = b_0$.
    Note that by polynomial multiplication $\varphi(p(x)q(x)) = a_0b_0 = \varphi(p(x))\varphi(q(x))$ and clearly $\varphi$ is a surjection.
    Also note that every $p(x) \in \ker(\phi)$ must have zero constant term and therefore is comprised of products and sums of $x$ so $\ker(\varphi) = (x)$.
    %Let $p(x) \in \ker(\varphi)$ and note that $p(x)$ must have a zero constant term and therefore is comprised of products and sums of $x$ thus 
    %$p(x) \in (x)$, so $\ker(\varphi) \subseteq (x)$. Let $p(x) \in (x)$ and suppose for the sake of contradiction that $p(x)$ has a nonzero constant term, $a_0$. By definition $p(x) - a_0 \in (x)$ and it must 
    %follow that $p(x) - (p(x) - a_0) = a_0 \in (x)$ which implies that $a_0(a_0)^{-1} = 0 \in (x)$. Thus $(x) = R[x]$ and therefore $(x)$ is not a prime ideal. Since all $p(x) \in (x)$ have zero constant term 
    %$p(x) \in \ker(\varphi)$ so  $(x) \subseteq \ker(\varphi) $ and $\ker(\varphi) = (x)$. 

    By The First Isomorphism Theorem for Rings we find that $R[x]/(x) \cong R$. By Proposition 13, supposing $(x)$ is a prime ideal in $R[x]$ we know that $R[x]/(x)$ is an integral domain, thus $R$ is an integral domain. 
    Since Proposition 13 is an if and only if statement if we suppose $R$ is an integral domain then we get that $R/(x)$ is an integral domain and therefore $(x)$ is a prime ideal. 

    Supposing $(x)$ is a maximal ideal we get that $R/(x)$ is a field by Proposition 12 and therefore $R$ is a field. 
    Again since Proposition 12 is an if and only if statement, supposing $R$ is a field we know that $R/(x)$ is a field and therefore 
    $(x)$ must be a maximal ideal. 

  \end{proof}
\end{exercise}
\vspace{.5in}

\begin{exercise}[\textbf{7.4.30}] Let $I$ be the ideal of the commutative ring $R$ and define, 
  \begin{equation*}
    \text{rad }I = \{r \in R | r^n \in I \text{ for some } n \in \ZZ^{+} \}
  \end{equation*} 
  called the radicals of $I$. Prove that  $\text{rad }I$ is an ideal containing $I$ and that $\text{rad }I/I$ is the nilradical of quotient ring $R/I$.
  \begin{proof} Suppose $I$ is an ideal of the commutative ring $R$ and $\text{rad }I$ is defined as above. Clearly $I \subseteq  \text{rad }I$, 
    let $i \in I$ and note that $i ^1 \in I$ so $i \in  \text{rad }I$. Let $a, b \in \text{rad }I$, by definition there
    exists some $i, j \in \ZZ^{+}$ such that $a^i,b^j \in I$. Note that by binomial expansion we know that
    \begin{equation*}
      (a - b)^{i + j} = \sum_{k = 0}^{i + j} {i + j \choose k}a^{i + j - k}(-b)^k.
    \end{equation*}
    When $k\geq j$, since $R$ is commutative we can factor a $b^j$ term and we conclude that every term in the sum where 
    $k\geq j$ is in $I$. When $k < j$, since $R$ is commutative we can factor an $a^i$ term and we conclude that every term in the sum where 
    $k < j$ is in $I$. Thus every term in the sum is in $I$ and $(a - b)^{i + j} \in I$. Thus by definition $(a - b) \in \text{rad }I$.

    Since $R$ is commutative $(ab)^{ij} = a^{ij}b^{ij} = a^{ij}(b^{j})^i \in I$ so $ab \in \text{rad }I$. So $\text{rad }I$ is a subring of $R$.

    Suppose $r \in  R$ and consider that since $R$ is commutative $(ar)^i = a^ir^i \in I$. Therefore $ar \in  \text{rad }I$, since $I$ is an ideal the same argument can be applied for 
    $ra \in  \text{rad }I$. Thus $ \text{rad }I$ is an ideal in $R$.  

    The nilradical of the quotient ring $R/I$ is the ideal containing all the nilpotent elements of $R/I$. By definition if an element $(x + I)\in R/I$ 
    is nilpotent then there exists some $n$ such that $(x + I)^n = x^n + I = I$ (well defined since $R$ is commutative, and recall that $I$ is the zero in $R/I$).
    Thus $x^n \in \text{rad }I$ and we conclude that $\text{rad }I/I$ is the nilradical of quotient ring $R/I$.
  \end{proof}
\end{exercise}
\vspace{.5in}



\begin{exercise}[\textbf{7.4.41}] A proper ideal $Q$ of the commutative ring $R$ is called primary if whenever 
  $ab \in Q$ and $a \not\in Q$ then $b^n  \in Q$ for some positive integer $n$. Establish the following facts about 
  primary ideals, 
  \begin{enumerate}
  \item[(a)] The primary ideals of $\ZZ$ are $0$ and $(p^n)$, where $p$ is a prime and $n$ is a positive integer.
  \begin{proof} Clearly $0$ is a primary ideal of $\ZZ$. Let $ab \in (p^n)$ then either $p\mid a$ or $p\mid b$ which implies either $a^n  \in (p^n)$ or $b^n  \in (p^n)$.
    Suppose for the sake of contradiction $(i)$ is a prime ideal where $i$ is not prime or prime power. Let $i = p_1^{\alpha_1}p_2^{\alpha_2} \dots p_n^{\alpha_n}$ where $p_k$ are prime and 
    note that by definition $(p_1^{\alpha_1}p_2^{\alpha_2} \dots p_{n - 1}^{\alpha_{n - 1}})(p_n^{\alpha_n}) = i \in (i)$ however clearly neither $p_n^{\alpha_n} \in (i)$ or $(p_1^{\alpha_1}p_2^{\alpha_2} \dots p_{n - 1}^{\alpha_{n - 1}}) \in (i)$
    so $(i)$ is not primary. 
  \end{proof}
\vspace{.15in}

\item[(b)] Every prime ideal of $R$ is a primary ideal. 
\begin{proof} Suppose $P$ is a prime ideal. By definition if $ab \in P$ then either $a \in P$ or $b \in P$. Without loss of generality suppose 
  $a \not\in P$ and $b \in P$, then clearly $b^n \in P$ and we find that $P$ is primary. 
\end{proof}
\vspace{.15in}

\item[(c)] An ideal $Q$ of $R$ is primary if and only if every zero divisor in $R/Q$ is a nilpotent element of $R/Q$. 
\begin{proof} Suppose $Q$, an ideal of $R$ is primary. Let $(i + Q) \in R/Q$ be a zero divisor, so $(j + Q) \in R/Q$ such that $(i + Q)(j + Q) = (ij + Q) = Q$ where $(i + Q),(j + Q) \neq Q$. Therefore 
  $ij \in Q$ and $j \not\in Q$ since $Q$ is primary, there exists some $i^n \in Q$ so we get $(i + Q)^n = (i^n + Q) = Q$. Thus $(i + Q)$ is nilpotent in $R/Q$. 

  Suppose every zero divisor in $R/Q$ is a nilpotent element of $R/Q$. Let $ab \in Q$ and $a \not\in Q$ note that $(ab + Q) = (a + Q)(b + Q) = (a + Q)(Q) = Q$. Thus 
  $(b + Q)$ is a zero divisor. Then there exists some $(b + Q)^n = (b^n + Q) = Q$ so $b^n \in Q$. Therefore $Q$ is primary. 
\end{proof}
\vspace{.15in}

\item[(d)] If $Q$ is a primary ideal then $\text{rad}(Q)$ is a prime ideal. 
\begin{proof} Suppose $Q$ is a primary ideal. Let $ab \in \text{rad}(Q)$. By definition $\text{rad}(Q)$ there exists some $n$ such that $(ab)^n \in Q$, 
  Since $R$ is commutative we know that $(a^n)(b^n) \in Q$, without loss of generality suppose $(a^n) \not\in Q$, since $Q$ is primary then there exists some 
  $b^{mn} \in Q$. Therefore $b \in \text{rad}(Q)$.
\end{proof}


  \end{enumerate}
\end{exercise}
\vspace{.5in}

\section*{\textbf{Section 9.1}}


\begin{exercise}[\textbf{9.1.1}] Let $p(x, y, z) = 2x^2y - 3xy^3z + 4y^2z^5$ and $q(x, y, z) = 7x^2 + 5x^2y^3z^4 + 3x^2z^3$ be polynomials 
  in $\ZZ[x, y, z]$.
  \begin{enumerate}
    \item[(a)] Write each of $p$ and $q$ as polynomials in $x$ with coefficients in $\ZZ[y, z]$
    \begin{proof} Writing $p$ and $q$ as polynomials in $x$ with coefficients in $\ZZ[y, z]$ we get, 
      \begin{equation*}
        p(x) = (4y^2z^5) - (3y^3z)x + (2y)x^2,
      \end{equation*}
      \begin{equation*}
        q(x) = (7 + 5y^3z^4 + 3z^3)x^2.
      \end{equation*}
    \end{proof}
\vspace{.15in}
      \item[(b)] Find the degree of each of $p$ and $q$.
      \begin{proof} By definition we know that the degree of a nonzero polynomial is the largest degree of any of it's monomial terms and the 
        degree of a monomial $x_1^{d_1}x_2^{d_2}x_3^{d_3}\dots x_n^{d_n}$ is simply $d_1 + d_2 + d_3 + \dots + d_n$. Thus $deg(p) = 7$ and $deg(q) = 9$.
      \end{proof}
\vspace{.15in}

      \item[(c)] Find the degree of $p$ and $q$ in each of the three variables $x, y,$ and $z$
      \begin{proof} As we have shown in part $a$, $deg_x(p) = 2$ and $deg_x(q) = 2$. Note that $p$ and $q$ as 
        polynomials in $y$ with coefficients in $\ZZ[x, z]$ we get, 
        \begin{equation*}
          p(y) = + (2x^2)y + (4z^5)y^2 - (3zx)y^3,
        \end{equation*}
        \begin{equation*}
          q(y) = (7x^2 + 3x^2z^3) + (5x^2z^4)y^3.
        \end{equation*}
        So $deg_y(p) = 3$ and $deg_y(q) = 3$. Note that $p$ and $q$ as 
        polynomials in $z$ with coefficients in $\ZZ[x, y]$ we get, 
        \begin{equation*}
          p(y) = 2x^2y - (3xy^3)z + (4y^2)z^5,
        \end{equation*}
        \begin{equation*}
          q(y) = 7x^2 + (3x^2)z^3 + (5x^2y^3)z^4.
        \end{equation*}
        So $deg_z(p) = 5$ and $deg_z(q) = 4$.
      \end{proof}
      \vspace{.15in}

      \item[(d)] Compute $pq$ and find the degree of $pq$ in each of the three variables $x, y,$ and $z$. 
      \begin{proof} Computing $pq$ we get, 
        \begin{align*}
          pq &= (2x^2y - 3xy^3z + 4y^2z^5)(7x^2 + 5x^2y^3z^4 + 3x^2z^3)\\
         &= (2x^2y - 3xy^3z + 4y^2z^5)7x^2 + (2x^2y - 3xy^3z + 4y^2z^5)5x^2y^3z^4 + (2x^2y - 3xy^3z + 4y^2z^5)3x^2z^3\\
         &= 2x^2y(7x^2) - 3xy^3z(7x^2) + 4y^2z^5(7x^2) + 2x^2y(5x^2y^3z^4) - 3xy^3z(5x^2y^3z^4) \\
         &+ 4y^2z^5(5x^2y^3z^4) + 2x^2y(3x^2z^3) - 3xy^3z(3x^2z^3) + 4y^2z^5(3x^2z^3)\\
         &= 14x^4y - 21x^3y^3z + 28y^2z^5x^2 + 10x^4y^4z^4 - 15x^3y^6z^5 + 20x^2y^5z^9 + 6x^4yz^3 - 9x^3y^3z^4 + 12y^2z^8x^2
        \end{align*}
        Note that $deg(pq) = deg(p) + deg(q) = 16$. Finding the degree of $pq$ in each variable we get $deg_x(pq) = 4$, $deg_y(pq) = 6$, and $deg_z(pq) = 9$.
        \end{proof}
        \vspace{.15in}

      \item[(e)] Write $pq$ as a polynomial in the variable $z$ with coefficients in $\ZZ[x, y]$
      \begin{proof} Computing $pq$ as a polynomial in the variable $z$ with coefficients in $\ZZ[x, y]$ we get, 
        \begin{equation*}
          pq(z) = 14x^4y - (21x^3y^3)z + (6x^4y)z^3 + (10x^4y^4- 9x^3y^3)z^4 + (28y^2x^2 - 15x^3y^6)z^5 + (12y^2x^2)z^8 + (20x^2y^5)z^9. 
        \end{equation*}
        \end{proof}
  \end{enumerate}
\end{exercise}
\vspace{.5in}



\begin{exercise}{\textbf{9.1.4}} Prove that the ideals $(x)$ and $(x, y)$ are prime ideals in $\QQ[x, y]$ but only the latter ideal is 
  a maximal ideal. 
  \begin{proof} Suppose $p, q \in \QQ[x, y]$ such that $pq \in (x)$. Writing $p$ and $q$ as polynomials in the variable $x$ with coefficients in $\QQ[y]$ 
    we get, 
    \begin{equation*}
      p(x) = a_0 + a_1x + a_2x^2 + \dots + a_nx^n
    \end{equation*} 
    \begin{equation*}
      q(x) = b_0 + b_1x + b_2x^2 + \dots + b_mx^m
    \end{equation*} 
    Now consider $pq$ and note that $a_0b_0$ is a term in $pq$ which is degree $0$ in the variable $x$, however since $pq \in (x)$ it 
    must be the case that $a_0b_0 = 0$. By Proposition 1 since $\QQ$ is an integral domain we know that $\QQ[y]$ is an integral domain so 
    it must follow that either $a_0 = 0$ or $b_0 = 0$. Without loss of generality suppose $a_0 = 0$, then 
    we can factor an $x$ out of $p$ and we conclude that $p \in (x)$.

    Suppose $p, q \in \QQ[x, y]$ such that $pq \in (x, y)$. Again it must follow that the constant terms in $a_0$ and $b_0$ in $p$ and $q$ must 
    have the property that $a_0b_0 = 0$. By Proposition 1 since $\QQ$ is an integral domain we know that $\QQ[x,y]$ is an integral domain so 
    it must follow that either $a_0 = 0$ or $b_0 = 0$. Without loss of generality suppose $a_0 = 0$, then every term $p \in (x, y)$. 

    Apply Exercise 7.4.7 to show that $(x, y)$ is maximal in $\QQ(x, y)$ by noting that $\QQ$ is a field.
  \end{proof}
\end{exercise}
\vspace{.5in}



\begin{exercise}{\textbf{9.1.9}} Prove that a polynomial ring in infinitely many variables with coefficients in any commutative ring contains ideals that are not finitely 
  generated.  
  \begin{proof}  Suppose $R$ is a commutative ring and $R[x_1, x_2 \dots]$ and for the sake of contradiction suppose $I = (x_1, x_2 \dots x_n)$ is an ideal. Let $x_m \in R[x_1, x_2 \dots]$ such that 
    $m > n$, and note that since $I$ is an ideal then $x_mI \in I$, however by definition clearly any monomial with $x_m$ is not in $I$. 

  \end{proof}
\end{exercise}
\vspace{.5in}







\end{document} 


